\section{Methodology}\label{sec:methodology}
This section first formulates the LiDAR NeRF-BA problem. We then describe the insights on volume sampling for LiDAR NeRF and the proposed LiDAR NeRF-BA method. 

\subsection{Problem Formulation: LiDAR NeRF-BA}
In general, a NeRF model can be formulated as: 
\begin{equation}\label{eq:NeRF_model}
	\begin{split}
		\hat{\sigma} = f\left(\gamma\left(T(\mathbf{x};\mathbf{p})\right);\Theta\right)
	\end{split}
\end{equation}
where $\hat{\sigma}$ is the estimated density. $T(\cdot)$ denotes the Euclidean transformation function that transforms a volume sample $\mathbf{x} \in \mathbb{R}^3$ by sensor pose $\mathbf{p} \subset SE(3)$. $\gamma(\cdot)$ is the positional encoding function, and $f(\cdot)$ is a multi-layer perceptron (MLP) with parameters $\Theta$, representing the NeRF model. 
Note that the viewing direction in the original five-dimensional NeRF input~\cite{mildenhall2020nerfOG} affects only view-dependent radiance, not density. As our model predicts density alone, volume samples remain three-dimensional although we continue to refer to the model as a NeRF throughout this paper.

For any LiDAR ray denoted as $\mathbf{r}(z)=\mathbf{o}+z\mathbf{u}$, where $\mathbf{o}$ and $\mathbf{u}$ are the ray's origin and direction, with $z \in \mathbb{R}_{>0}$ being the range along the LiDAR ray, an estimated range $\hat{D}$ from the NeRF model can be derived by the alpha compositing function of $N$ volume samples \cite{mildenhall2020nerfOG}:
\begin{equation}\label{eq:alpha_compositing_part1}
	\hat{D}(\mathbf{r}) = \sum_{i=1}^{N} \hat{h}(z_i) z_i 
\end{equation}
\begin{equation}\label{eq:termination_distribution}
	\hat{h}(z_i) =
	\exp\left(- \sum_{j=1}^{i-1} \hat{\sigma}_j \delta_j \right) 
	\left(1 - \exp\left(-\hat{\sigma}_i \delta_i\right)\right) 
\end{equation}
where $\delta_i$ is the distance between neighboring samples, and $\hat{h}(z)$ is the predicted termination distribution of the ray.

Given $S$ number of LiDAR scans of point clouds, and the initial estimates of their sensor poses $P = \{\mathbf{p}_s\}_{s=0}^S$, LiDAR NeRF-BA can be formulated as the following optimization problem:

\begin{equation}\label{eq:NeLD_BA_problem_def}
	P^*, \Theta^* = \operatorname*{arg\,min}_{P,\,\Theta}\;
	\mathcal{L}\!\left(\hat{D}(P, \Theta),\; D\right)
\end{equation}
where $D$ is the measured range from LiDAR sensors, and each $\mathbf{p}_s$ is parameterized as the Lie algebra $se(3)$ of $SE(3)$. The optimization goal is to find the set of optimal sensor poses $P^*$ and MLP parameters $\Theta^*$ that minimize a chosen loss function $\mathcal{L}$ (see Section \ref{sec:loss}). 

\subsection{Volume Sampling for LiDAR NeRF-BA}\label{sec:volume_sampling}

%--- Image Here---
\begin{figure*}[h!]
	\centering
	\includegraphics[width=0.9\textwidth]{images/rgb_nerf_ba_vs_lidar_nerf_ba_5.png}
	\caption{\textbf{Comparison of RGB NeRF-BA and LiDAR NeRF-BA in Volume Sampling along Depth/Range.}
		(a) In RGB NeRF-BA, depth is only \textit{implicitly} constrained via colors and multi-view geometry (orange rays), thus volume samples are loosely constrained. In contrast, LiDAR provides range measurements which require \textit{explicit} volume sampling to constrain. 
		(b) Sampling around the sharp transition of $\sigma(z)$ (pink region) is critical to produce a difference in the ray termination distribution $h(z)$ and $h(z+\epsilon)$ after alpha blending, which provides non-vanishing $\frac{\partial \hat{h}(z_i)}{\partial T(\mathbf{x}_j; \mathbf{p})}$ for pose optimization.
		(c) NeRF \cite{mildenhall2020nerfOG} allocates fine samples \textit{around} coarse samples, which can miss the sharp transition of $\sigma(z)$ (highlighted pink region). We instead sample \textit{prior to} coarse samples, ensuring a dense sampling around the high-probability spike in $\sigma$, akin to a step function.}
	\label{fig:rgb_nerf_ba_vs_lidar_nerf_ba}
\end{figure*}

Compared to RGB NeRF-BA which uses photometric
loss for optimizing camera poses~\cite{lin2021barf,wang2022nerfmm,chen2023l2gnerf}, LiDAR NeRF-BA should make full use of the \textit{range loss} due to the strong geometric cue from LiDAR rays. While RGB NeRF-BA relies on multi-view geometry to optimize poses along depth, LiDAR NeRF-BA should use \textit{range}, as shown in Figure~\ref{fig:rgb_nerf_ba_vs_lidar_nerf_ba} (a).

This motivates us to consider the gradient of the estimated range $\hat{D}$ with respect to the sensor pose $\mathbf{p}$:

\begin{equation}\label{eq:range_gradient}
\frac{\partial \hat{D}}{\partial \mathbf{p}}
=
\sum_{i=1}^{N} z_i
\sum_{j=1}^{i}
\frac{\partial \hat{h}(z_i)}{\partial T(\mathbf{x}_j; \mathbf{p})}
\frac{\partial T(\mathbf{x}_j; \mathbf{p})}{\partial \mathbf{p}}
\end{equation}

We have found that the density of volume sampling, i.e., the number of samples per unit length along the ray, is essential for ensuring the $\frac{\partial \hat{h}(z_i)}{\partial T(\mathbf{x}_j; \mathbf{p})}$ terms do not vanish, as illustrated in Figure \ref{fig:rgb_nerf_ba_vs_lidar_nerf_ba}(b). We address this through two modifications to volume sampling. 

First, note that the vanilla hierarchical sampling strategy~\cite{mildenhall2020nerfOG} is sub-optimal for providing gradients for optimising $\mathbf{p}$. In the vanilla approach, fine samples were placed in a bin \textit{around} coarse samples, guided by weights derived from the coarse termination distribution, as shown in Figure \ref{fig:rgb_nerf_ba_vs_lidar_nerf_ba}(c). Notably, the sum in equation \ref{eq:range_gradient} is typically dominated by the first sample after the $\sigma$ spike, and samples after it contribute little to the sum due to alpha blending. Thus, we propose to place fine samples at volumes \textit{before} coarse samples, as shown in Figure~\ref{fig:rgb_nerf_ba_vs_lidar_nerf_ba} (c). This fully utilizes the fine samples and significantly eliminates the estimated range ambiguity, as shown in Figure~\ref{fig:ray_termination_dist_comparison} with an actual LiDAR ray. 

Second, we notice that LiDAR rays often travel tens of meters into empty space, leading to sparse volume sampling. We therefore bound each ray's sampling interval to the normalized cube $C = [-1,1]^3$. Specifically, we first find the intersections of each ray with the cube $C$ using the AABB slab method~\cite{williams2005raybox}, which gives the near intersection range $c_n$ and far intersection range $c_f$. We then clamp $c_n$ and $c_f$ with chosen lower and upper bounds, $z_{low}$ and $z_{up}$, to get the sampling interval $[z_n, z_f]$ for the ray:

\begin{equation}\label{eq:cube_bound_conditions}
z_n = \max(z_\text{low},\, c_n), \qquad
z_f = \min(z_\text{up},\, c_f),
\end{equation}
with a fallback to $[z_n, z_f] = [z_\text{low}, z_\text{up}]$ when the ray does not intersect $C$. Note that this bounding is only possible for point clouds, where the 3D positions of points are provided, unlike images with a potentially unbounded scene. It is most effective for rays whose origin is close to the scene boundary and points away from $[0,0,0]$, and significantly improves the overall volume sampling density.

\begin{figure}
	\centering
	\includegraphics[width=0.47\textwidth]{images/ray_term_dist_comparison_3.png}
	\caption{\textbf{Ray Termination Distributions of a LiDAR Ray in canteen\_day sequence} Our sampling method better approximates the baseline distribution constructed with 8192 uniform samples when shifted ($\epsilon=0.01$), which is critical for providing accurate gradients in LiDAR NeRF-BA.}
	\label{fig:ray_termination_dist_comparison}
	\vspace{-15pt}
\end{figure}

\subsection{On Fourier Positional Encoding}
Positional encoding is key to detailed scene reconstruction in both RGB NeRF and LiDAR NeRF~\cite{mildenhall2020nerfOG}. The Fourier positional encoding function $\gamma(\cdot)$ for feature $x$ is \cite{tancik2020fourierfeatures}:
\begin{align}\label{eq:fourier_positional_encoding_0}
	\gamma(x) & = \left[x, \gamma_0(x), \gamma_1(x), \dots, \gamma_{L-1}(x)\right] \in \mathbb{R}^{1 + 2L}
\end{align}
where $\gamma_k(x) = \left[\cos\left(2^k \pi x\right), \sin\left(2^k \pi x\right)\right] \in \mathbb{R}^{2}$.

There is, however, a $2^k\pi$ amplification in the encoding gradient that makes it difficult to optimize $\mathbf{p}$ and may drive $\mathbf{p}$ away from the global minimum. This is particularly important for the common outdoor use case of LiDAR scenes that uses high-frequency encoding for reconstructing scene details~\cite{martinbrualla2021nerfinwild}. Inspired by the successful use of \textit{surrogate gradient} to optimize spiking neural networks~\cite{Zenke2021surrogate}, the following surrogate gradient for Fourier positional encoding is proposed to handle this problem:

\begin{align}\label{eq:surrogate_fourier_gradient}
	{\frac{\partial \gamma_k^\text{smooth}(x)}{\partial x}}
	= \frac{1}{2^k\pi} {\frac{\partial \gamma_k(x)}{\partial x}}  = \begin{bmatrix}
		- \sin(2^k \pi x) , \cos(2^k \pi x)
	\end{bmatrix}
\end{align}

The normal positional encoding function in equation~\ref{eq:fourier_positional_encoding_0} was used in the forward pass, and the gradient function in equation~\ref{eq:surrogate_fourier_gradient} was used in the backward pass to avoid gradient amplification. Clearly, $\gamma^\text{smooth}_k(\cdot)$ preserves the gradient direction of $\gamma_k(\cdot)$ while removing the amplification component. This allows for using high-frequency Fourier features with no amplification on gradients, ensuring $\mathbf{p}$ stays at the global minimum even for large $L$.
Additionally, we followed BARF~\cite{lin2021barf} in gradually activating Fourier features, and optimization of sensor poses starts only when the $k^*$th frequency for position features is activated. This allows a coarse reconstruction to first be reached by the NeRF model, and the sensor pose can then be more stably optimized towards the global minimum.

% --- Raw Map Image result [Overview | zoomin]
\begin{figure*}
	\vspace*{5pt}
	\centering
	\renewcommand{\arraystretch}{0.4}
	\begin{tabular}{@{}c@{\hskip 2pt}c@{\hskip 2pt}c@{\hskip 2pt}c@{\hskip 2pt}c@{\hskip 2pt}c@{}}
		\rotatebox{90}{\parbox[c]{2cm}{\centering\textbf{\small quad\_easy (NC)}}} &
			\includegraphics[width=0.2\textwidth]{images/overview_map/quad_easy.png} &
			\includegraphics[width=0.18\textwidth]{images/raw_map/quad_easy/fastlio_initial.png} &
			\includegraphics[width=0.18\textwidth]{images/raw_map/quad_easy/hba.png}     &
			\includegraphics[width=0.18\textwidth]{images/raw_map/quad_easy/balm.png}    &
			\includegraphics[width=0.18\textwidth]{images/raw_map/quad_easy/neld_ba.png} \\

		\rotatebox{90}{\parbox[c]{2cm}{\centering\textbf{\small quad\_hard (NC)}}} &
			\includegraphics[width=0.2\textwidth]{images/overview_map/quad_hard.png} &
			\includegraphics[width=0.18\textwidth]{images/raw_map/quad_hard/fastlio_initial.png} &
			\includegraphics[width=0.18\textwidth]{images/raw_map/quad_hard/hba.png}     &
			\includegraphics[width=0.18\textwidth]{images/raw_map/quad_hard/balm.png}    &
			\includegraphics[width=0.18\textwidth]{images/raw_map/quad_hard/neld_ba.png} \\

		\rotatebox{90}{\parbox[c]{2cm}{\centering\textbf{\small math\_hard (NC)}}} &
			\includegraphics[width=0.2\textwidth]{images/overview_map/math_hard.png} &
			\includegraphics[width=0.18\textwidth]{images/raw_map/math_hard/fastlio_initial.png} &
			\includegraphics[width=0.18\textwidth]{images/raw_map/math_hard/hba.png}&
			\includegraphics[width=0.18\textwidth]{images/raw_map/math_hard/balm.png}    &
			\includegraphics[width=0.18\textwidth]{images/raw_map/math_hard/neld_ba.png} \\

		\rotatebox{90}{\parbox[c]{2cm}{\centering\textbf{\small canteen\_day (FP)}}} &
			\includegraphics[width=0.2\textwidth]{images/overview_map/canteen_day.png} &
			\includegraphics[width=0.18\textwidth]{images/raw_map/canteen_day/fastlio_initial.png} &
			\includegraphics[width=0.18\textwidth]{images/raw_map/canteen_day/hba.png}     &
			\includegraphics[width=0.18\textwidth]{images/raw_map/canteen_day/balm.png}    &
			\includegraphics[width=0.18\textwidth]{images/raw_map/canteen_day/neld_ba.png} \\

		\rotatebox{90}{\parbox[c]{2cm}{\centering\textbf{\small garden\_day (FP)}}} &
			\includegraphics[width=0.20\textwidth]{images/overview_map/garden_day.png} &
			\includegraphics[width=0.18\textwidth]{images/raw_map/garden_day/fastlio_initial.png} &
			\includegraphics[width=0.18\textwidth]{images/raw_map/garden_day/hba.png}     &
			\includegraphics[width=0.18\textwidth]{images/raw_map/garden_day/balm.png}    &
			\includegraphics[width=0.18\textwidth]{images/raw_map/garden_day/neld_ba.png} \\[6pt]
		&
        	\small Ground Truth Map &
			\small FAST-LIO2~\cite{xu2021fastlio2} (initial) & 
			\small HBA~\cite{liu2022HBA} & 
			\small BALM ~\cite{liu2021balm} & 
			\small NeLD-BA (ours)\\
	\end{tabular}
	\caption{\textbf{Qualitative Results of Raw Point Cloud Maps for Newer College (NC) and FusionPortable (FP) Sequences.}}
	\label{fig:raw_map_results}
\end{figure*}

\subsection{Loss Functions}\label{sec:loss}
To train the LiDAR NeRF-BA model, the L1 loss was used on both the coarse and fine estimated ranges:
% --- equation here ---
\begin{align}\label{eq:depth_intensity_loss}
	\mathcal{L}_\text{range}  = \sum_\mathbf{r}\left\| \hat{D}(\mathbf{r}) - D(\mathbf{r}) \right\|_1
\end{align}
The distribution loss in DS-NeRF~\cite{deng2024dsnerf} is also used for ray termination distributions, which is the KL divergence loss for comparing the estimated termination distribution $\hat{h}(z)$ against a targeted termination distribution $h^*(z)$:
% --- equation here ---
\begin{align}\label{eq:termination_distribution_loss}
	\mathcal{L}_\text{termination} = \sum_\mathbf{r}\sum_{n=1}^N h^*(z_n) \cdot ln\left(\frac{h^*(z_n)}{\hat{h}(z_n)}\right)
\end{align}
While DS-NeRF modeled the target termination distribution as a Gaussian distribution with a small variance, we leverage the fact that range is provided in LiDAR measurements and propose a pseudo Dirac delta function as the targeted termination distribution. This effectively enforces that the ray termination probability in the first sample after ray termination is 1:
% --- equation here ---
\begin{align}
	\label{eq:pseudo_dirac_delta}
	h^*(z_n) =
	\begin{cases}
		1, & \text{if } n = \underset{j \in \mathbb{N},\; z_j - D > 0}{\arg\min} { \left\| z_j - D \right\| } \\
		0, & \text{otherwise}
	\end{cases}
\end{align}

Our final loss function $\mathcal{L}$ is:
% --- equation here ---
\begin{align}
\mathcal{L} 
&= \lambda_\text{coarse} \cdot 
   \bigl(\lambda_d \mathcal{L}_\text{range}^c 
       + \lambda_h \mathcal{L}_\text{termination}^c \bigr) \notag \\
&\quad + (\lambda_d \mathcal{L}_\text{range}^f
       + \lambda_h \mathcal{L}_\text{termination}^f)
\label{eq:final_loss}
\end{align}
where superscripts $c$ and $f$ denote loss functions for coarse and fine estimations, and $\lambda$s are hyperparameters to be tuned.

\subsection{Mapping}
A rendered 3D map can be built by casting rays from the optimized sensor poses with the original ray directions in LiDAR frames. Using the same hierarchical sampling strategy mentioned above, the termination distances $z_\text{map}$ of rays are determined to be at the position of the highest termination probability per unit length:
\begin{align}
	\label{eq:mapping_with_termination_density}
	i^* = \arg\max_{i} \left( \frac{\hat{h}(z_i)}{\delta_i} \right),
	\qquad
	z_{\text{map}} = z_{i^*}
\end{align}
Finally, we combine points at $\mathbf{r}(z_\text{map})$ from all frames to create a rendered map.

\subsection{Implementation Details}\label{par:implementation_detail}
The same hyperparameters were used for all sequences. In our implementation, each scan was voxelized with a grid size of 0.05 m. % all frames and sensor positions were uniformly normalized together into the $[-1,1]^3$. 
Our model architecture follows NeRF in the Wild~\cite{martinbrualla2021nerfinwild}, with a hidden layer size of 512, positional encoding of $L=15$ for position features and $L=4$ for direction features. All models were trained for 160k iterations. Fourier features were gradually activated using BARF's method~\cite{lin2021barf} from iteration 0 to 120k. Sensor poses were unfrozen when $k^*=6$ frequency for position features was activated.  

The number of coarse and fine samples was increased from 32 to 64 and from 16 to 32 from iteration 0 to iteration 120k. Learning rates were exponentially decayed from $5\times10^{-5}$ to $2\times10^{-5}$. % $2\times10^{-4}$ to $2\times10^{-5}$ (rotation), and $5\times10^{-4}$ to $5\times10^{-5}$ (position). 
The default PyTorch~\cite{paszke2019pytorch} implementation of the Adam optimizer~\cite{kingma2017adam} was used for training. 

Hyperparameters were obtained through grid search with $\lambda_\text{coarse}=0.1$, $\lambda_\text{fine}=1$, $\lambda_\text{d}=1$, $\lambda_\text{i}=1$ and $\lambda_\text{h}=1$. Rendered maps were generated using 128 coarse samples and 128 fine samples with our proposed volume sampling strategy. All models took approximately 1.5 hours to train on a single NVIDIA 4060 GPU.